\Needspace{8\baselineskip}
\Article{203}
\textsc{Problem.} \textit{If the forms \(\Phi,\varphi\) are equivalent, give all transformations of one into the other.}

\textit{Solution.} When these forms are equivalent in only one way (\textit{i.e.}, either only properly or only improperly), find by Article~196 one transformation of \(\varphi\) into \(\Phi\), say \(\alpha,\beta,\gamma,\delta\). Clearly no transformations other than those similar to this one can exist. But when \(\varphi,\Phi\) are both properly and improperly equivalent, find two dissimilar transformations, \textit{i.e.}, one proper and one improper, say \(\alpha,\beta,\gamma,\delta\) and \(\alpha',\beta',\gamma',\delta'\). Every other transformation will be similar to one or the other. So let the form \(\varphi\) be \((a,b,c)\), its determinant \(=D\), the greatest common divisor of \(a,2b,c\) be \(m\) (as throughout the preceding discussion), and let \(t,u\) range over all numbers satisfying \(tt-Duu=mm\). In the first case all transformations of \(\varphi\) into \(\Phi\) are contained in the first of the following formulas, I; in the second case, either in the first, I, or in the second, II.
\[
\begin{array}{@{}lll@{}}
\mathrm{I}\;\ldotp\,\ldotp\,\ldotp\,\ldotp\,\ldotp\,\ldotp\,\ldotp&\frac1m(\alpha t-(\alpha b+\gamma c)u),&\frac1m(\beta t-(\beta b+\delta c)u)\\[1ex]
&\frac1m(\gamma t+(\alpha a+\gamma b)u),&\frac1m(\delta t+(\beta a+\delta b)u)\\[1.5ex]
\mathrm{II}\;\ldotp\,\ldotp\,\ldotp\,\ldotp\,\ldotp\,\ldotp\,\ldotp&\frac1m(\alpha't-(\alpha'b+\gamma'c)u),&\frac1m(\beta't-(\beta'b+\delta'c)u)\\[1ex]
&\frac1m(\gamma't+(\alpha'a+\gamma'b)u),&\frac1m(\delta't+(\beta'a+\delta'b)u)
\end{array}
\]
\SourceLeaf{202}
\textit{Example.} - We seek all transformations of $(129,92,65)$ into $(42,59,81)$. In Article~195 we found that these are equivalent only improperly, and in the next article we obtained the improper transformation $-47,-56,73,87$ of the former into the latter. So all transformations of $(129,92,65)$ into $(42,59,81)$ are given by the formula
\[
-(47t+421u),\quad -(56t+503u).\quad 73t+653u,\quad 87t+780u
\]
where $t$, $u$ range over all numbers satisfying $tt-79uu=1$. These are given by
\[
\begin{aligned}
\pm t&=\tfrac12\bigl((80+9\sqrt{79})^e+(80-9\sqrt{79})^e\bigr)\\
\pm u&=\tfrac{1}{2\sqrt{79}}\bigl((80+9\sqrt{79})^e-(80-9\sqrt{79})^e\bigr)
\end{aligned}
\]
where $e$ takes all nonnegative integers.

\Needspace{8\baselineskip}
\Article{204}
Clearly, the general formula giving all transformations becomes \textit{simpler} as the initial transformation from which it is derived becomes simpler. Since we can choose which transformation to start from, we can often make the general formula simpler by deriving a simpler transformation from the formula first found, giving $t$, $u$ particular values, and then forming another formula from it. For \textit{example}, putting $t=80$, $u=-9$ in the formula found in the preceding article's example gives a simpler transformation than the one we started from, namely $29,47,-37,-60$. From this we derive the general formula $29t-263u$, $47t-424u$, $-37t+337u$, $-60t+543u$. So once we have obtained a general formula by the preceding rules, we can try giving $t$, $u$ particular values $\pm t',\pm u';\ \pm t'',\pm u''$ etc., to see whether we get a simpler transformation than the one from which the formula was derived. If so, we can derive a simpler formula from that transformation. --- Some freedom remains in judging simplicity. If it were worth the effort, we could put this on a fixed footing, and also assign \textit{bounds} in the sequence $t',u';\ t'',u''$ etc. beyond which the transformations become continually less simple. We would then need to try only those within the bounds, without going further. But with the methods we have given, the simplest transformation usually appears either at
\SourceLeaf{203}
once or on taking $t$, $u$ to have the values $\pm t',\pm u'$. So to keep this account brief, we leave out this investigation.

\Needspace{8\baselineskip}
\Article{205}
\textsc{Problem.} \textit{Find all representations of a given number $M$ by a given formula $axx+2bxy+cyy$, with positive nonsquare determinant $=D$.}

\textit{Solution.} First we observe that finding representations in values of $x$, $y$ that are not relatively prime can be reduced to the case of relatively prime values of the indeterminates, in exactly the same way as above for forms with negative determinant (Article~181). So it would be unnecessary to repeat this here. For representations in relatively prime values of $x,y$ to be possible, $D$ must be a quadratic residue of $M$. If all values of the expression $\sqrt D\;(\mathrm{mod.}\,M)$ are $N,-N,N',-N',N'',-N''$ etc. (which we may take so that none is $>\tfrac12M$), every representation of $M$ by the given form belongs to one of these values. We must therefore first find these values, and then investigate the representations belonging to each in turn. There are no representations belonging to $N$ unless the forms $(a,b,c)$ and $(M,N,\frac{NN-D}{M})$ are properly equivalent. If they are, find any proper transformation of the former into the latter, say $\alpha,\beta,\gamma,\delta$. We then have the representation of $M$ by $(a,b,c)$ belonging to $N$, namely $x=\alpha$, $y=\gamma$, and all representations belonging to this value are given by
\[
x=\tfrac1m\bigl(\alpha t-(\alpha b+\gamma c)u\bigr),\qquad
y=\tfrac1m\bigl(\gamma t+(\alpha a+\gamma b)u\bigr)
\]
where $m$ denotes the greatest common divisor of $a,2b,c$, and $t$, $u$ range over all numbers satisfying $tt-Duu=mm$. --- Clearly this general formula becomes simpler as the transformation $\alpha,\beta,\gamma,\delta$ from which it is derived becomes simpler. So it is useful first to find the simplest transformation of $(a,b,c)$ into $(M,N,\frac{NN-D}{M})$, as in the preceding article, and derive the formula from it. --- In exactly the same way, general formulas give the representations belonging to the remaining values $-N,N',-N'$ etc., if there are any.

\textit{Example.} We seek all representations of $585$ by the formula
\SourceLeaf{204}
$42xx+62xy+21yy$. For representations in values of $x$, $y$ that are not relatively prime, it is immediately clear that the only possible ones have greatest common divisor of $x$, $y$ equal to $3$, since $585$ is divisible by only one square, $9$. So once we have found all representations of $\frac{585}{9}$, \textit{i.e.}, $65$, by the form $42x'x'+62x'y'+21y'y'$, in which $x'$ is relatively prime to $y'$, all representations of $585$ by $42xx+62xy+21yy$ in which $x$ is not relatively prime to $y$ follow by putting $x=3x'$, $y=3y'$. The values of $\sqrt{79}\;(\mathrm{mod.}\,65)$ are $\pm12,\pm27$. A representation of $65$ belonging to $-12$ is $x'=2$, $y'=-1$. So all representations of $65$ belonging to this value are given by $x'=2t-41u$, $y'=-t+53u$, and all representations of $585$ arising \textit{from these} by $x=6t-123u$, $y=-3t+159u$. In the same way we find the general formula giving all representations of $65$ belonging to $+12$, namely $x'=22t-199u$, $y'=-23t+211u$; and the formula containing all representations of $585$ arising from these, $x=66t-597u$, $y'=-69t+633u$. For the values $+27$ and $-27$, however, there is no representation of $65$. --- To find the representations of $585$ in relatively prime values of $x$, $y$, we must first find the values of $\sqrt{79}\;(\mathrm{mod.}\,585)$, which are $\pm77,\pm103,\pm157,\pm248$. No representation is found to belong to $\pm77,\pm103,\pm248$. But to $-157$ belongs the representation $x=3$, $y=1$, from which we derive the general formula giving all representations belonging to this value, $x=3t-114u$, $y=t+157u$. Similarly, we find the representation belonging to $+157$, namely $x=83$, $y=-87$, and the formula containing all similar representations, $x=83t-746u$, $y=-87t+789u$. So we have four general formulas containing all representations of $585$ by $42xx+62xy+21yy$:
\[
\begin{aligned}
x&=6t-123u&\qquad y&=-3t+159u\\
x&=66t-597u& y&=-69t+633u\\
x&=3t-114u& y&=t+157u\\
x&=83t-746u& y&=-87t+789u
\end{aligned}
\]
where $t,u$ range over all integers satisfying $tt-79uu=1$.

To keep this account brief, we will not dwell on special applications of the preceding investigations into forms with positive nonsquare deter\SourceLeaf{205}minant. Anyone can work these out without difficulty, along the lines of Articles~176, 182. We turn at once to forms with positive square determinant, the only ones left to consider.
